% The GRASS Temporal Framework: Object oriented code design with examples.
% Build with make (see README.md); the examples run as tests in
% temporal/tests/temporal_framework_design_test.py.
\documentclass[10pt,a4paper]{article}

% A plain info dictionary, so that the test reads /Keywords from the bytes of
% the PDF, and no date or random identifier, so that the same sources give
% the same PDF.
\pdfobjcompresslevel=0
\pdfinfoomitdate=1
\pdftrailerid{}
\pdfsuppressptexinfo=-1

\usepackage[T1]{fontenc}
\usepackage[utf8]{inputenc}
\usepackage{lmodern}
\usepackage{microtype}
\usepackage[a4paper,margin=2.2cm]{geometry}
\usepackage{xcolor}
\usepackage{booktabs}
\usepackage{caption}
\usepackage{float}
\usepackage{placeins}
\usepackage{listings}
\usepackage[numbers]{natbib}
\usepackage[simplified]{pgf-umlcd}
\usetikzlibrary{fit,calc}
% PDFDocEncoding writes the digest of the keywords as plain ASCII.
\usepackage[pdfencoding=pdfdoc]{hyperref}

\InputIfFileExists{digest.tex}{}{\newcommand{\sourcedigest}{unknown}}

\hypersetup{
  pdftitle={The GRASS Temporal Framework: Object oriented code design with
    examples},
  pdfauthor={S\"oren Gebbert},
  pdfkeywords={\sourcedigest},
  colorlinks=true,
  linkcolor=black,
  citecolor=black,
  urlcolor=blue!50!black,
}

\renewcommand{\umldrawcolor}{black}
\renewcommand{\umlfillcolor}{gray!8}

% A figure file uses plain underscores in names of classes and members.
% The underscore of the roman font is one and a half letters wide, so that
% _PlanarExtentMethods looked like __PlanarExtentMethods, a mangled name in
% Python. In the typewriter font, as in the text and the listings, it is as
% wide as a letter.
\newcommand{\umlfigure}[1]{%
  \begingroup\catcode`\_=12 \ttfamily\input{figures/#1}\endgroup}

% Code that a reader copies from the PDF must run: monospaced, straight
% quotes, spaces kept, no line breaks or numbers added. A space of a listing
% is otherwise only a gap between glyphs: mutool draw -F txt and pdftotext
% -raw gave the code without indentation (measured 2026-11-06, MuPDF 1.23.10,
% poppler 24.02). Here each space is the visible-space glyph, drawn in
% invisible text mode and mapped to U+0020, so the text layer holds one real
% space per space of the file at the width of a character. Measured the same
% day: mutool draw -F txt then keeps the indentation; pdftotext of poppler
% 24.02 still drops the spaces at the start of a line, with -raw and without
% options; -layout, which places text by its position, gave 3 of the 8
% examples exactly.
\input{glyphtounicode}
\pdfgentounicode=1
\pdfglyphtounicode{uni2423}{0020}
\makeatletter
\def\lst@visiblespace{\pdfliteral direct{3 Tr}\char32\pdfliteral direct{0 Tr}}
\makeatother
\lstdefinestyle{example}{
  language=Python,
  basicstyle=\ttfamily\small,
  columns=fullflexible,
  showspaces=true,
  upquote=true,
  showstringspaces=false,
  breaklines=false,
  frame=lines,
  keywordstyle=\color{blue!50!black},
  commentstyle=\color{gray!70!black},
  stringstyle=\color{green!35!black},
}
\lstdefinestyle{output}{
  basicstyle=\ttfamily\small,
  columns=fullflexible,
  showspaces=true,
  upquote=true,
  breaklines=false,
  frame=single,
  rulecolor=\color{gray},
}
\renewcommand{\lstlistingname}{Listing}

\newcommand{\code}[1]{\texttt{#1}}
% The path of an example file in the GRASS repository.
\newcommand{\examplefile}[1]{%
  \code{python/grass/docs/src/temporal\_framework\_design/examples/#1}}
% Long names in \texttt do not break; let the lines around them stretch.
\setlength{\emergencystretch}{3em}

\title{The GRASS Temporal Framework:\\
  Object oriented code design with examples}
\author{S\"oren Gebbert}
\date{GRASS 8.6, 6 November 2026}

\begin{document}

\maketitle

\begin{abstract}
\noindent This document is a revised version of the supplementary material of
S\"oren Gebbert and Edzer Pebesma: The GRASS GIS temporal framework.
International Journal of Geographical Information Science 31(7), 1273--1292,
2017, doi:10.1080/13658816.2017.1306862 \citep{gebbert2017}. It describes the
code as it is in GRASS 8.6. The version of 2017, by S\"oren Gebbert and Edzer
Pebesma, lies unchanged next to this document as
\code{Temporal-Framework-API-Description.pdf}; it shows the code of 2017.
\end{abstract}

\tableofcontents

\section{Introduction}

The temporal framework of GRASS manages time series of raster, 3D raster and
vector maps \citep{gebbert2014,gebbert2017}. A map gets a time stamp, and
maps with time stamps are registered in space time datasets: space time
raster datasets (STRDS), space time 3D raster datasets (STR3DS) and space time
vector datasets (STVDS). The time stamps and the spatial extents of maps and
datasets are kept in a temporal database, by default an SQLite database in the
mapset (\code{tgis/sqlite.db}). The temporal tools (\code{t.*}) are written in
Python on top of the Python package \code{grass.temporal}, except
\code{t.connect}, which is written in C.

This document describes the classes of \code{grass.temporal}
(section~\ref{sec:design}), the database layer below them
(section~\ref{sec:database}) and the algebra core, which computes the
temporal algebra, the aggregations, the map calculations and the sampling of
the tools (section~\ref{sec:algebra}).
Section~\ref{sec:examples} shows Python examples: the registration of maps,
temporal aggregation, the same aggregation with the algebra core, sampling at
vector points, the spatio-temporal intersection of trajectories and a
benchmark.

Every example is a Python file of the GRASS repository, in the folder
\code{examples} of the source of this document,
\code{python/grass/docs/src/temporal\_framework\_design}; the caption of each
listing gives the path of its file. The
examples run in the North Carolina sample dataset \code{nc\_spm\_08\_grass7}
\citep{neteler2008}, available from \url{https://grass.osgeo.org/download/data/},
in a new mapset \code{temporal\_design}. They make their own data from maps of
the dataset and never write into its mapset \code{PERMANENT}. To run them, save
the files and start them in the order of section~\ref{sec:examples}, the first
one with \code{-c}, which creates the mapset:

\begin{lstlisting}[style=output]
grass -c ~/grassdata/nc_spm_08_grass7/temporal_design --exec python3 register.py
grass ~/grassdata/nc_spm_08_grass7/temporal_design --exec python3 aggregate.py
\end{lstlisting}

\noindent The tests of GRASS run each example in such a mapset and compare its
output with the output that this document shows; they check also that each
class and member that a figure draws exists in the code.

\section{Object oriented code design}
\label{sec:design}

Maps and space time datasets are Python objects whose content is stored in
the temporal database. Each object keeps its content in rows, one row per
table: the base information, the absolute or relative time, the spatial
extent, the metadata and a register. The figures of this section show the
classes in UML, simplified: they show the members that a user of the class
needs most, not all of them.

\subsection{Rows of the temporal database}

Figure~\ref{fig:rows} shows the classes of the rows. \code{DictSQLSerializer}
holds the content of a row in the dictionary \code{D} and turns it into the
SELECT, INSERT and UPDATE statements of a table and back.
\code{SQLDatabaseInterface} names the table and the identifier of the row and
runs the statements: \code{select()}, \code{insert()}, \code{update()},
\code{update\_all()}, \code{delete()} and \code{is\_in\_db()}. The identifier of a
map is \code{name@mapset}, of a vector map with a layer \code{name:layer@mapset},
because a time stamp can be given to each layer of a vector map.

\code{Row} lists its columns once, as \code{Column} objects in the class
attribute \code{columns}; the getter, the setter and the property of each
column are generated when the class is created. The rows of a map or dataset
are \code{BaseRow} (identifier, name, mapset, creator, creation time, temporal
type), \code{RegisterRow} (the datasets a map is registered in),
\code{AbsoluteTimeRow} or \code{RelativeTimeRow} (the time stamp),
\code{SpatialExtentRow} (the spatial extent) and a metadata row. One class
serves all six types of maps and datasets: a \code{Kind} (\code{kinds.py})
tells the row its table and which columns it has. The modification time and
the semantic type of \code{BaseRow}, and the granularity and the temporal type
of the maps (\code{map\_time}) of the time rows, are columns of space time
datasets only.

\code{TemporalExtent} computes the temporal relations between two time
stamps (the relations of \citet{allen1983} and their names in GRASS, such as
\code{starts()}, \code{during()}, \code{follows()} or \code{overlaps()},
together with \code{temporal\_relation()}, which names the relation), and the
intersection and the unions of two time stamps.
\code{SpatialExtent} does the same for spatial extents in two and three
dimensions; the methods with \code{\_2d} use north, south, east and west only.

\begin{figure}[p]
\centering
\umlfigure{rows}
\caption{The row classes. Every row has the identifier of its map or dataset
and the name of its table; a \code{Kind} chooses table and columns, so one
class serves raster, 3D raster and vector maps and their space time
datasets.}
\label{fig:rows}
\end{figure}

The metadata differ by type (figure~\ref{fig:metadata}). The six metadata
classes are generated from tables of columns in \code{metadata.py};
\code{MetadataRow(kind)} returns an object of the class of the kind, so that a
raster map has no number of points and a vector map no semantic label. The
metadata of a space time dataset hold the extremes of the metadata of its maps:
the smallest and largest minimum (\code{min\_min}, \code{min\_max}), the
smallest and largest maximum (\code{max\_min}, \code{max\_max}) and the
smallest and largest resolution in each direction (\code{nsres\_min} to
\code{tbres\_max}), because maps of different resolution can be registered in
one dataset. The mixins \code{\_SpaceTimeMetadata} and
\code{\_SemanticLabelsMetadata} add what the datasets share.

\begin{figure}[ht]
\centering
\umlfigure{metadata}
\caption{The metadata classes, generated from tables of columns. A list of
columns that ends in ``\ldots'' is shortened.}
\label{fig:metadata}
\end{figure}

\subsection{Abstract datasets and the topology}

Maps and space time datasets share an abstract base class,
\code{AbstractDataset} (figure~\ref{fig:datasets}). It reads the time stamp and
the spatial extent of an object, computes temporal relations and
intersections, and selects, inserts, updates and deletes the object in the
temporal database. The spatial methods are abstract: two mixins implement them
for two- and three-dimensional extents (figures~\ref{fig:maps}
and~\ref{fig:stds}).

\code{AbstractDataset} inherits from \code{TopologyLinks}, which keeps the
links of an object to the objects it is temporally and spatially related to:
one dictionary per kind of relation, with a list of objects for each relation,
and the temporal successor and predecessor (\code{next()}, \code{prev()}).
The \code{SpatioTemporalTopologyBuilder} sets these links. It builds the
topology of one list of objects, or between two lists, with the R*-tree of the
GRASS vector library, and it links the objects of each list in the order of
their start times by \code{next()} and \code{prev()}. An object that starts or finishes another
one is also entered as during it, and an object that is started or finished by
another one also contains it, so \code{get\_contains()} of a month lists every
day of the month, the first and the last day included.

\code{AbstractMapDataset} sets the time stamp of a map, buffers its time stamp
and its spatial extent, and keeps the register of the datasets a map is
registered in (\code{get\_registered\_stds()}, \code{add\_stds\_to\_register()},
\code{remove\_stds\_from\_register()}); \code{delete()} removes a map from every
dataset it is registered in. \code{MapDataset} implements the methods that
read a map: \code{load()} asks the C library for the metadata of the map, in a
separate process (\code{c\_libraries\_interface.py}), so that an error of the C
library does not end the caller.

\code{AbstractSpaceTimeDataset} registers and unregisters maps,
updates the dataset from its maps (\code{update\_from\_registered\_maps()}),
computes the granularity and the temporal topology of its maps, and shifts and
snaps their time stamps. Shifting moves all time stamps by a granule, such as
one year, two weeks or five minutes. Snapping sets the end of each map to the
start of the next one, except for maps that start together; a last map without
an end time gets one granule of the granularity of the dataset as its length,
a last map with an end time keeps it. The
register of the maps of a dataset is a table of its own, which
\code{StdsMapRegister} creates, fills and reads.

\begin{figure}[p]
\centering
\umlfigure{datasets}
\caption{The abstract dataset classes. Abstract methods are in italics. The
topology builder links the objects of one or two lists; the register of the
maps of a space time dataset is a table of its own.}
\label{fig:datasets}
\end{figure}

\subsection{Map and space time dataset classes}

The classes of the three map types and of the three space time dataset types
(figures~\ref{fig:maps} and~\ref{fig:stds}) combine \code{MapDataset} or
\code{StdsDataset} with a mixin for the spatial extent, \code{\_PlanarExtentMethods}
for two dimensions or \code{\_VolumeExtentMethods} for three. A class sets its
\code{Kind}; a map class also copies what the C library reports about a map of
its type into the metadata. The rows are made by \code{reset()}, which the
constructor calls with the identifier. Of the classes of these figures, the
six classes of the types are exported from \code{grass.temporal};
\code{MapDataset}, \code{StdsDataset} and the mixins are not.

\begin{figure}[ht]
\centering
\umlfigure{maps}
\caption{The map classes and the rows of a map.}
\label{fig:maps}
\end{figure}

\begin{figure}[ht]
\centering
\umlfigure{stds}
\caption{The space time dataset classes and the rows of a dataset. A dataset
has the register of its maps instead of a register of datasets.}
\label{fig:stds}
\end{figure}

% No figure of the sections before may come out inside the listings below.
\FloatBarrier
\subsection{Registered maps as lists of objects}

Three methods of a space time dataset return its maps as objects. The
objects carry the identifier, the time stamp and the spatial extent of a map,
which is what the topology needs; \code{select()} reads the rest.
Figures~\ref{fig:as-objects} to~\ref{fig:by-granularity} show the result
for a dataset of three maps of absolute time with gaps between them: A from
20 to 22 January 2010, B from 24 to 27 January and C from 28 to 29 January.
The granularity of the dataset is one day. The arrows show the order of the
list. They are not links between the objects: these methods set no
\code{next()} or \code{prev()}. Listing~\ref{lst:map-lists} creates the dataset
and prints the three lists.

\begin{figure}[ht]
\centering
\umlfigure{maps_as_objects}
\caption{The maps of the dataset ordered by start time. The arrows show the
order of the list; they are not links between the objects, the method sets no
\code{next()} or \code{prev()}.}
\label{fig:as-objects}
\end{figure}

\begin{figure}[ht]
\centering
\umlfigure{maps_with_gaps}
\caption{The maps with a gap object for each gap. A gap is an object of the
class of the maps, here \code{RasterDataset}, with the identifier
\code{None} and the time of the gap. The arrows show the order of the list;
they are not links between the objects, the method sets no \code{next()} or
\code{prev()}.}
\label{fig:with-gaps}
\end{figure}

\begin{figure}[ht]
\centering
\umlfigure{maps_by_granularity}
\caption{One list per granule of the granularity. Each list holds a new object
for each map in the granule, with the identifier and the spatial extent of the
map and the time of the granule, or a gap object. The arrows show the order of
the lists; they are not links between the objects, the method sets no
\code{next()} or \code{prev()}. The method needs a valid temporal topology:
maps must not overlap or contain each other.}
\label{fig:by-granularity}
\end{figure}

% A figure that LaTeX puts at the top of a page must not stand between two
% pages of a listing, which a reader copies as one.
\FloatBarrier

\lstinputlisting[style=example,caption={The dataset of
figures~\ref{fig:as-objects} to~\ref{fig:by-granularity}
\\ \examplefile{map\_lists.py}},label=lst:map-lists]{examples/map_lists.py}

\lstinputlisting[style=output,caption={Output of
\code{map\_lists.py}}]{examples/map_lists.txt}

\section{The database layer}
\label{sec:database}

The database layer consists of five modules:

\begin{description}
\item[\code{state}] \code{init()} finds out everything about the session
  once: the database backend, the database, the mapset, the messenger and the
  C library interface. All of it lives in one object, \code{STATE}; the
  functions \code{get\_*()} read it, and only \code{init()} and
  \code{set\_raise\_on\_error()} write it. Unless it is called with
  \code{skip\_db\_init}, \code{init()} creates the temporal database if it
  does not exist and stops at a database of a version that the framework does
  not read; \code{t.upgrade} upgrades an older one. Every script calls
  \code{init()} before it uses the framework.
\item[\code{db}] \code{SQLDatabaseInterfaceConnection} holds the connections
  to the temporal databases of the mapsets in the search path, one
  \code{DBConnection} per database; a mapset outside the search path joins
  when a dataset or map named with it is first read. The backend is SQLite
  (\code{sqlite3}) or PostgreSQL (\code{psycopg2}). The functions
  \code{sql\_identifier()} and \code{sql\_statement()} check every table,
  view and column name that goes into the text of a statement; values are
  passed as parameters or quoted by \code{sql\_literal()}.
\item[\code{schema}] creates the tables and views of the temporal database
  from the SQL templates, reads its version from its metadata table and
  upgrades it (\code{t.upgrade}).
\item[\code{kinds}] names the six types in one place: the names a user may
  type, the tables and views of each type, the prefix of the functions of the
  C library interface, the type of \code{g.remove} and the module that
  computes a map of the type from an expression.
\item[\code{rows}] the row classes of figure~\ref{fig:rows}.
\end{description}

\noindent A method that reads or writes the database takes an optional
\code{dbif}. Without it, the method opens a connection, runs its statements and
closes the connection again; a script that runs many statements opens one
connection and passes it on (listing~\ref{lst:register}).
\code{register\_maps\_in\_space\_time\_dataset()} always opens a connection
of its own; its argument \code{dbif} is deprecated.

\section{The algebra core and the classic algebra}
\label{sec:algebra}

The algebra core is the package \code{grass.temporal.algebra}. It computes the
expressions of \code{t.rast.algebra}, \code{t.rast3d.algebra}, \code{t.vect.algebra} and
\code{t.select}, the aggregations of \code{t.rast.aggregate} and
\code{t.rast.aggregate.ds}, the map calculations of \code{t.rast.mapcalc} and
\code{t.rast3d.mapcalc} and the sampling of \code{t.sample} and
\code{t.vect.observe.strds} (figure~\ref{fig:algebra}). A series is an
immutable, time ordered sequence of slots. A slot has a time stamp
(\code{Extent}), optionally a spatial extent (\code{Space}), and a content tree
that says what the map of the slot contains: an existing map (\code{MapRef}),
a constant (\code{Const}), a function of the raster algebra (\code{Call}) or a
reduction of several maps (\code{SeriesReduce}, \code{SeriesAggregate}).
Building a series computes no map.

The operations are functions from series to series: \code{select()} keeps the
slots that are related to slots of another series, \code{pairs()} combines
related slots, \code{gather()} reduces the slots related to each slot of a
target, \code{apply()} and \code{where()} combine series slot by slot. A
\code{Relation}, such as \code{contains} or \code{during|starts}, says which
slots are related; the core asks the topology builder of
\code{grass.temporal}. \code{source()} reads a space time dataset as a series
and \code{granules()} makes empty slots on a regular grid of time. Only
\code{materialize()} computes maps: it turns the content trees into calls of
\code{r.mapcalc}, \code{r3.mapcalc}, \code{r.series}, \code{g.copy},
\code{v.overlay}, \code{v.patch} or \code{v.buffer} and registers the maps in
a new space time dataset. \code{plan()} returns the same calls as a
\code{ProcessChain} without running them.

The parser classes of the tools, such as \code{TemporalRasterAlgebraParser},
read an expression with a grammar and a translator, which turns the syntax
tree into calls of the operations. A \code{Profile} collects the rules of a
language: \code{CORE} for the core, \code{CLASSIC} for the classic algebra.

The classic algebra is the language of the temporal algebra as the manuals and
the publications define it. It is selected with the flag \code{-c} of the
algebra tools and \code{t.select} or with \code{GRASS\_TEMPORAL\_ALGEBRA=classic}, and it is
computed on the core by the profile \code{CLASSIC}. The implementation that is
older than the core stays in \code{grass.temporal.classic} as the reference
implementation, because published results rely on it; the tests compare the
classic profile with it. It changes only together with documented corrections
of the classic algebra.

Listing~\ref{lst:algebra} computes the monthly means of
listing~\ref{lst:aggregate} with the core.

\begin{figure}[ht]
\centering
\umlfigure{algebra}
\caption{The algebra core. The box \code{algebra} is the package
\code{grass.temporal.algebra} with some of its functions; the other boxes are
classes. Series, slots and content trees are immutable values.}
\label{fig:algebra}
\end{figure}

% Nor inside the listings of the examples.
\FloatBarrier
\section{Code examples}
\label{sec:examples}

The examples run in the mapset \code{temporal\_design} of the North Carolina
sample dataset, in the order of this section; an example names the examples
it needs in its docstring. Run times depend on the computer; the output does
not, except for the times that the benchmark writes to standard error. The
tests compare the output of each example with the output shown here.

\subsection{Registration of raster maps}

Listing~\ref{lst:register} computes 366 maps of the daily mean air
temperature of 2020 with one call of \code{r.mapcalc}. The values are
generated, not measured, and stored as whole hundredths of a degree Celsius,
as the gridded E-OBS data that the version of 2017 used. Each map is loaded as
a \code{RasterDataset}, gets its day as time stamp, is inserted into the
temporal database and registered in a new STRDS. One database connection
serves all statements.

\lstinputlisting[style=example,caption={Registration of daily temperature
maps \\ \examplefile{register.py}},label=lst:register]{examples/register.py}

\lstinputlisting[style=output,caption={Output of
\code{register.py}}]{examples/register.txt}

\subsection{Temporal aggregation}

Listing~\ref{lst:aggregate} aggregates the daily maps to monthly means in
degree Celsius. A granule is a \code{RasterDataset} without a map: it has only
a time stamp, here one month. The topology builder relates the twelve granules
to the daily maps, \code{r.series} averages the days that a granule contains,
and \code{r.mapcalc} divides by 100. Each result map gets the time stamp of
its granule and is registered in a new STRDS. The objects of
\code{get\_registered\_maps\_as\_objects()} carry the identifier and the
extents only, so the script reads the metadata with \code{select()} before
it prints the range of each month.

\lstinputlisting[style=example,caption={Aggregation of daily to monthly
temperature \\ \examplefile{aggregate.py}},label=lst:aggregate]{examples/aggregate.py}

\lstinputlisting[style=output,caption={Output of
\code{aggregate.py}}]{examples/aggregate.txt}

\subsection{Aggregation with the algebra core}

Listing~\ref{lst:algebra} computes the same monthly means with the algebra
core: \code{granules()} gives the months, \code{gather()} averages the days
that each month contains with \code{r.series}, and the division by 100 is a
slot by slot operation. \code{materialize()} computes the maps and registers
them. The values are the same as those of listing~\ref{lst:aggregate}.

\lstinputlisting[style=example,caption={Monthly means with the algebra core
\\ \examplefile{algebra.py}},label=lst:algebra]{examples/algebra.py}

\lstinputlisting[style=output,caption={Output of
\code{algebra.py}}]{examples/algebra.txt}

\subsection{Sampling a raster time series at vector points}

Listing~\ref{lst:sample} samples the monthly means at three climate stations of
the dataset: on the coast, in the Piedmont and in the mountains. It copies the
stations into the mapset with \code{v.extract}, reads their coordinates with
pygrass and reads the raster cell under each station directly from each
monthly map.

\lstinputlisting[style=example,caption={Sampling at climate stations
\\ \examplefile{sample.py}},label=lst:sample]{examples/sample.py}

\lstinputlisting[style=output,caption={Output of
\code{sample.py}}]{examples/sample.txt}

\subsection{Spatio-temporal intersection of trajectories}

The version of 2017 intersected the GPS tracks of two zebras. Here the tracks
are generated from the dataset: two field crews sample the five stream gauges
of Crabtree Creek in Raleigh (\code{usgsgages}) on one morning in opposite
order, so that they meet on the way (listing~\ref{lst:tracks}). The positions
are used directly, without vector maps: each becomes a \code{VectorDataset}
without a map, whose time stamp is the time of the position and whose spatial
extent is the position (listing~\ref{lst:trajectories}). Snapping makes each
position last until the next one and the last one for the granularity of the
positions, five minutes; buffering widens a position by half the distance to
the next position, at least by 50~m. The topology builder relates the
positions of the two crews in two dimensions; a pair that is related in space
and in time is a meeting, printed with the time and the area of the
intersection.

\lstinputlisting[style=example,caption={The tracks of two field crews
\\ \examplefile{tracks.py}},label=lst:tracks]{examples/tracks.py}

\lstinputlisting[style=output,caption={Output of
\code{tracks.py}}]{examples/tracks.txt}

\lstinputlisting[style=example,caption={Where and when the crews met
\\ \examplefile{trajectories.py}},label=lst:trajectories]{examples/trajectories.py}

\lstinputlisting[style=output,caption={Output of
\code{trajectories.py}}]{examples/trajectories.txt}

\subsection{Benchmark}

Listing~\ref{lst:benchmark} measures typical operations: the registration of
maps, the selection of maps by SQL as rows and as objects, and the temporal
topology with the count of the relations. The number of maps is the first
argument; the output shown is that of 50 maps, as the tests run it. The
counts go to standard output, the times to standard error.

\lstinputlisting[style=example,caption={Benchmark
\\ \examplefile{benchmark.py}},label=lst:benchmark]{examples/benchmark.py}

\lstinputlisting[style=output,caption={Output of \code{benchmark.py 50}}]%
{examples/benchmark.txt}

Table~\ref{tab:benchmark} gives the times of \code{benchmark.py 1000},
measured on 6 November 2026 on an Intel Core Ultra 9 285T with 24 cores,
GRASS 8.6.0dev, Python 3.12 and SQLite 3.45: the median of three runs, each in
a new mapset of the dataset. They show the order of magnitude on one computer,
not a property of the framework.

\begin{table}[ht]
\centering
\caption{Times of \code{benchmark.py 1000}, median of three runs}
\label{tab:benchmark}
\begin{tabular}{lr}
\toprule
Operation & Time [s] \\
\midrule
Registration & 19.4 \\
SQL selection & 0.011 \\
Selection as objects & 0.026 \\
Topology and counting & 0.026 \\
\bottomrule
\end{tabular}
\end{table}

\bibliographystyle{plainnat}
\bibliography{references}

\end{document}
